\chapter{Motivation}

The motivation for TrustBridge arises from three observable gaps in the current crowdfunding landscape.

\textbf{Gap 1 — Contributor Information Asymmetry.} Contributors to crowdfunding campaigns face a fundamental information asymmetry: they must evaluate campaign quality based entirely on creator-supplied material, without any independent, data-driven assessment. Research by Elitzur et al. (2024) and Ardakani et al. (2024) demonstrates that ML models trained on historical crowdfunding data can predict campaign success with meaningful accuracy using only information available at launch time. Yet no mainstream crowdfunding platform exposes such a signal to contributors. TrustBridge closes this gap by computing and displaying an ML-derived success probability and an AI-generated risk explanation on every campaign page.

\textbf{Gap 2 — Lack of Programmable Post-Funding Accountability.} In conventional crowdfunding, once a campaign is successfully funded, the release of money is governed by platform policy — not by code. A creator who fails to deliver has, in practice, already received the funds. Smart-contract-based crowdfunding research (reviewed in Section 5) demonstrates that blockchain escrow can enforce predefined release conditions without relying on platform intermediaries. TrustBridge implements a four-tranche milestone escrow in Solidity: an initial 20\% development tranche is released upon goal attainment, and the remaining 80\% is released in three subsequent tranches only after an authorised verifier approves the corresponding milestone evidence. If a milestone is rejected, the contract executes a predefined refund rule, returning the unspent escrow to contributors.

\textbf{Gap 3 — No Integrated Architecture Combining All Three Layers.} Individually, ML crowdfunding prediction, blockchain escrow, and explainable AI have each been studied in the literature. What has not been evaluated is an integrated system in which these three layers operate together: ML and Agentic AI at the assessment stage, smart-contract escrow at the funding stage, and AI-plus-human verification at the milestone stage. TrustBridge is motivated by the hypothesis that this integration produces a qualitatively different — and more trustworthy — crowdfunding experience than any single layer can provide alone.

Beyond the academic motivation, the team is personally motivated by the practical problem of trust in early-stage technology funding in India, where contributors increasingly encounter projects that are either misrepresented or fail to deliver, with no recourse available.
